\documentclass[a4paper,11pt]{article}

% Formatting
\usepackage[hmargin=2.5cm,top=2cm,bottom=3cm]{geometry}
\usepackage[T1]{fontenc}
\usepackage{microtype}
\usepackage{hyperref}
\usepackage{enumitem}

% Math 
\usepackage{amsmath,amssymb,amsthm}
\usepackage{mathrsfs}
\usepackage{mathtools}
\usepackage[capitalise]{cleveref}
\usepackage{bm}

% Operators
\DeclareMathOperator{\MO}{MO}
\DeclareMathOperator{\Span}{Span}

% Discussion
\usepackage{xcolor}
\newcommand{\question}[1]{\textcolor{blue}{[#1]}}

% Title
\title{Adaptive basis set generation methods}
\author{}
\date{\today}

\begin{document}

\maketitle

\section{Introduction}

The Density-Based Basis-Set Correction (DBBSC) method \cite{diata_phd}
is an electronic structure theory method based on the 
one-electron density. This method introduces a functional that corrects 
electron correlation and allows to achieve convergence to the
complete basis set limit. The high accuracy of DBCBS makes it interesting for 
use in solving electronic structure theory problems on quantum simulators, 
for future quantum computer benchmarking purposes. We study numerical tools that
make the implementation of this method possible on quantum simulators with a 
number of available qubits up to 28. This corresponds to a number of doubly 
occupied molecular orbitals \question{spin-orbitals?} equal to 14.

Existing molecular orbital discretisation methods include wavelets
\cite{genovese2022flexibilities} and linear combinations of atomic orbitals,
for instance Gaussian-type orbitals. We focus on the
latter. Existing methods for compressing molecular orbitals, such as the active space
method \cite{wang2023basis}, do not easily allow to reduce the number of MOs to
minimal sizes while respecting chemical accuracy. Other existing methods based on 
the Cholesky decomposition \cite{hummel2017low} of the Hamiltonian operator
employ low rank approximations to improve storage requirements and operation count. 
Such class of methods is based on matrix or tensor compression 
from a linear algebra point of view. Here, we investigate alternative 
density-based solutions, that achieve low computational and storage cost by 
approximating the electronic density function using smaller basis sets. 
The advantage of this approach is that the resulting smaller basis set 
can be used for multiple purposes. Namely, it can directly be used as a 
discretisation basis for minimizing any given energy functional, since it is 
Hamiltonian-independent. 

The main motivation behind our approach is the presence of linear dependencies,
which are intrinsic to any given atomic orbital basis set of Gaussian-type orbitals. 
Inspired by density fitting methods \cite{julien_df}, our strategy
aims at detecting linear dependencies. This technique produces a smaller basis set using a 
density-based criterion, which is the approximation of a given target electron density under 
any given norm. In the present work, we propose new adaptive basis set 
generation methods based on density fitting. We divide the computation time into
an \textit{offline phase}, performed on a classical computer, and an
\textit{online phase}, performed on a quantum simulator. The calculation of the
target electron density, used as a reference, and the adapted basis generation
are both performed during the offline phase.

The present manuscript is organized as follows. The first parts describes the
mathematical background of our adapted basis set generation methods and the
second part contains numerical results. The performance of 
our method is numerically assessed on small molecules.

\section{Theory}

We briefly describe the theory of the DBBSC method, 
originally proposed in \cite{giner2018curing}. 
Let us consider a system of $N$ electrons and the electronic Hamiltonian
\[
    H = T + W_{ee} + V_{ne},
\]
where $T$ is the kinetic operator, $W_{ee}$ is the
electron repulsion operator, given by $\sum_{i<j}^Nr_{ij}^{-1}$, and $V_{ne}$ 
is the nuclei-electron attraction operator. The exact electronic ground state
energy problem reads
\begin{equation}
    \label{eq:pb_exact}
    E_0 = \min_{\Psi\in W}\,\langle \Psi,H\Psi\rangle,
\end{equation}
where the exact wave function lives in the space defined as
\[
    W := \{\Psi\in\wedge^N H^1(\mathbb R^3,\mathbb C)\colon \langle
    \Psi,\Psi\rangle=1\}.
\]
We discretise problem of \cref{eq:pb_exact} by introducing a given finite basis
set, denoted by $B\subset H^1(\mathbb R^3,\mathbb C)$, and the discretised wave
function space
\[
    W^B := \{\Psi\in\wedge^N \Span(B)\colon \langle\Psi,\Psi\rangle=1\}.
\]
The discretised ground state problem then reads
\begin{equation}
    \label{eq:pb_discrete}
    E_*^B = \min_{\Psi\in W^B}\,\langle \Psi, H\Psi\rangle.
\end{equation}
The limitation of estimating the exact ground state energy by solving 
the discretised problem of \cref{eq:pb_discrete} is that the 
error on the energy converges as $\mathcal O(|B|^{-1})$, due to the cusp of the
exact electron wave function. In order to improve the error convergence, the DBBSC method 
proposes to approximate the ground state energy by solving the following discretised problem
\begin{equation}
    \label{eq:pb_cbs}
    E_0^B = \min_{\Psi\in W^B}\, \left(\langle\Psi, H\Psi\rangle +
    \bar E^B[n_{\Psi}]\right),
\end{equation}
with $\bar E^B[n_{\Psi}]$ the basis set correction functional, we refer 
to the original work for its explicit definition.  
The one-electron density $n_{\Psi}$, for any given wave 
function $\Psi\in W^B$, lives in the space
\[
    D^B := \{n\in D\colon \exists \Psi\in W^B, \quad n_{\Psi}=n\},
\]
with 
\[
    D= \{n\in L^1(\mathbb R^3)\colon n\geq 0, \quad \int_{\mathbb R^3} 
    n=N, \quad \sqrt{n}\in H^1(\mathbb R^3)\}.
\]
The problem of \cref{eq:pb_cbs} is solved self-consistently 
\cite{giner2021self}.

\subsection{Slater determinants}

In Hartree-Fock and post Hartree-Fock methods, a special form of wave function
is considered, known as a \textit{Slater determinant}, defined as
\begin{equation}
    \label{eq:slater}
    \Psi(\bm r_1,\ldots,\bm r_N) = \frac{1}{\sqrt{N!}}
    \det\left(\phi_i(\bm r_j)\right)_{1\leq i,j\leq N}
\end{equation}
where each $\phi_i$ is a normalised one-electron function of $H^1(\mathbb
R^3,\mathbb C)$ referred to as Molecular Orbital (MO). It can be easily verified
that $\Psi$ in the form of \cref{eq:slater} is indeed an element of $W$.
For any given Slater determinant $\Psi\in W$, the one-electron density is 
given in terms of the Molecular Orbitals (MOs) as
\begin{equation}
    \label{eq:den_mo}
    n_{\Psi} = \sum_{i=1}^N |\phi_i|^2.
\end{equation}

\subsection{Approximating the density}

Under the Linear Combination of Atomic Orbitals (LCAO) 
approximation, the discretisation basis $B$ is denoted by
$B=\{\chi_{\mu}\}_{1\leq \mu\leq N_b}$, where each $\chi_{\mu}$ is an atomic
orbital basis function. Such discretisation yields an approximation of the 
exact MOs on the basis of AOs. The problem unknown is the 
coefficient matrix $\bm C\in\mathbb R^{N\times N_b}$ in the
decomposition
\[
    \forall 1\leq i\leq N, \quad \phi_i = \sum_{\mu=1}^{N_b} \bm C_{i\mu}
    \chi_{\mu}.
\]
Substituting this MO expression in the one-electron density of \cref{eq:den_mo},
one obtains the following expansion of the density on the AOs:
\begin{equation}
    \label{eq:n_prod}
    n_{\Psi} = \sum_{\mu=1}^{N_b}\sum_{\nu=1}^{N_b} D_{\mu\nu} \chi_\mu\chi_\nu,
\end{equation}
where $\bm D = \bm C^\top \bm C$ is the so-called \textit{density matrix}, of
dimension $N_b\times N_b$.

\subsection{Problem formulation}

The problem consists in constructing a subset $I\subseteq\{1,\ldots,N_b\}$ of
any given size $M$ such that if the restricted basis 
$B_I:=\{\chi_{\mu}\}_{\mu \in I}$ is used to discretise the problem of
\cref{eq:pb_cbs}, i.e. if the minimisation space is $W^{B_I}$ instead of $W^B$, 
then the accuracy loss on the estimated ground state energy with respect to 
using the full basis $B$, i.e. the error
\[
    |E_0^B - E_0^{B_I}|,
\]
remains small.

Our main idea for solving this problem can be described as follows. 
Let us denote by $\Psi_0\in W^B$ the minimiser of problem of \cref{eq:pb_cbs} 
and by $n_{\Psi_0}\in D^B$ the associated electron density. Given a target size
$M$, we seek the optimal subset of $B$ of the form $B_I$ with
$|I|=M$, obtained by solving the following optimisation problem
\[
    \min_{\substack{I\subseteq\{1,\ldots,N_b\} \\ |I|=M}}\,
    \min_{m\in D^{B_I}}\,\|n_{\Psi_0} - m\|,
\]
where $\|\cdot\|$ is any given norm. Note that the optimisation criterion 
we used is the orthogonal projection error (also known as best approximation
error) with respect to the reference electron density. 
In other words, the problem boils down to selecting the index subset of size
$M$ such that the best approximation error is minimised.

Due to the expression of the density given by \cref{eq:n_prod}, any function
$m\in D^{B_I}$ is of the form
\[
    m = \sum_{\mu\in I}\sum_{\nu\in I}C_{\mu\nu} \chi_{\mu}\chi_{\nu}.
\]

\subsection{Algorithms}

To summarize, the main ingredients given as an input to our algorithm is 
\begin{itemize}
    \item given molecular geometry,
    \item atomic orbital basis set,
    \item a starting density matrix.
\end{itemize}
We studied three different choices of density matrices: Hartree-Fock (HF),
Configuration Interaction Singles and Doubles (CISD) and
Configuration Interaction Perturbatively Selected Iteratively (CIPSI).

\subsubsection{Assembly of the Gram matrix}

The first step is the assembly of the Gram matrix (i.e. the overlap matrix) of
the orbital products, given a user-defined inner product
$\langle\cdot,\cdot\rangle$. Let us denote by $\bm T$ the tensor of dimension
$N_b\times N_b\times N_b\times N_b$, with entries 
defined as
\[
    T_{\mu\nu\kappa\lambda} = D_{\mu\nu}D_{\kappa\lambda} 
    \langle \chi_{\mu}\chi_{\nu},\chi_{\kappa}\chi_{\lambda}\rangle.
\]
The Gram matrix $\bm G$, of dimension $N_b^2\times N_b^2$ is then formed by
folding pairwise the first two and the last two dimensions of the tensor 
$\bm T$. Then we also regroup by orbital types.
Another option for initialising the Gram matrix consists in $\bm
R^\top \bm G\bm R$, where $\bm R$ is the atomic restriction matrix that discards
orbital products centered on midbonds.

\subsubsection{Selecting orbital products}

We seek a selection of orbital products of target size $M_*\ll N_b^2$ 
that best approximate the one-electron density. This selection can be performed
using a criterion based on minimizing the orthogonal projection error with
respect to the full space spanned by orbital products, in which the density 
belongs to. Orbital products are discarded using a pivoted Cholesky 
decomposition performed on the Gram matrix $\bm G$, 
also known as Gram-Schmidt procedure with pivoting. The orthonormalization
procedure terminates when the orthogonal projection error is smaller than a
given threshold. Note that the entries of the Gram matrix are weighted by the
density matrix. Therefore, at the first iteration of the pivoted Cholesky 
decomposition, the selected pivot is influenced by the value of the density
matrix. 


The pivoted Cholesky decomposition under a given tolerance value
$\varepsilon>0$, performed on the matrix $\bm G$, yields an approximate rank
$M\leq N_b$ of $\bm G$ and a selection of $M$ rows (and columns) that 
approximately span the full row space under the tolerance value. We propose to
use the atomic orbitals corresponding to the selected indices to build the set
of adaptive basis functions. This strategy leads to our algorithm named ABS-0,
for Adaptive Basis Set generation.

Additionally, it may be of interest to impose maximal orbital type constraints
to the adaptive basis set, in favour of computational efficiency within the
integral evaluation procedure. The preliminary treatment of the Gram matrix $\bm G$ 
that enables this constraint can be presented as follows. The main idea is to
construct an intermediary atomic orbital basis respecting orbital type
constraints. We propose two possible
methods. For each $1\leq \mu\leq N_b$, if the atomic orbital $\chi_{\mu}$ has orbital 
type strictly higher than $d$-type, then
\begin{enumerate}[label=(\alph*)]
    \item ABS-1: discard $\chi_{\mu}$.
    \item ABS-2: for each orbital type $t\in\{s,p,d\}$, create a new atomic orbital
    having the orbital exponent and atomic center of $\chi_{\mu}$, but orbital
    type $t$.
\end{enumerate}
It is assumed that all atomic basis functions are normalized, so that the
diagonal of $\bm G$ consists of one entries.


\section{Numerical experiments}

\subsection{Implementation details}

The quantum simulator uses the Hamiltonian in the second quantisation
formulation, that reads
\begin{equation}
    \label{eq:H_sq}
    \hat H = \sum_{\sigma}\sum_{ij}^{N_{\MO}} h_{ij} \hat a^\dagger_{i\sigma} \hat
    a_{j\sigma} + \sum_{ijkl}^{N_{\MO}} h_{ijkl} \hat a^\dagger_{i\sigma} \hat
    a^\dagger_{k\sigma} \hat a_{j\sigma}\hat a_{l\sigma},
\end{equation}
where $N_{\MO}$ is the number of basis functions used to discretise the
molecular orbitals. Note that our method allows to reduce the number of MOs in this
expression.

\subsection{Results}

The first experiment allows to compare the results obtained from a DFT/B3LYP
calculation on the water molecule at equilibrium geometry, using
state-of-the-art auxiliary basis sets (RI and JKFIT) and our adaptive basis
sets, obtained with methods ABS-$x$ for $x=0,1,2$. 
The aim of this comparison is to study which method achieves better
accuracy with less basis functions. The approximation error is measured on
ground state energies using the full atomic orbital basis set as a reference
calculation. We are also interested in assessing the convergence of the
approximation error of our method for varying tolerance values in the Cholesky
decomposition.

\section{Perspectives}

Our method can be combined with orbital exponent and contraction coefficient
optimization methods, namely with the one proposed in \cite{wang2023basis}. A
possible direction that allows to combine both methods is to first apply our
method for basis set size reduction, then apply their method on our resulting 
adapted basis set for further optimization, in the goal of increasing the 
accuracy using a fixed number of MOs.

\bibliography{references}
\bibliographystyle{plain}

\end{document}



